\section{Related Work and Positioning}

\paragraph{EM as an optimization object.}
The expectation-maximization algorithm and the missing-information principle
provide the classical reference point for this work \citep{dempster1977maximum}.
The free-energy view of EM makes the coordinate-ascent structure explicit
\citep{neal1998view}. EM has also been analyzed as a preconditioned gradient
method \citep{xu1996convergence}, which is the closest prior framing to our
claim that EM conditioning is structural. The uniform $1/2$ curvature bound we
use follows the multinomial-logistic Hessian bound of \citet{bohning1992multinomial}
and the Popoviciu variance inequality.

\paragraph{Classical EM failure modes.}
Classical EM is cleanest when the expected complete-data objective separates
over component-private parameters. Shared parameters, non-separable M-steps,
singular components, hard assignments, and diffuse responsibilities are known
failure modes. This paper identifies their neural counterparts inside the chain
rule.

\paragraph{Generative versus discriminative training.}
The empirical distinction between generative and discriminative training is
well known \citep{ng2002discriminative}. We do not claim merely that the
objectives differ. We identify a mechanism inside one architecture: a curvature
dichotomy, measured gradient dominance, near-orthogonality of the CE and EM
paths, and causal restoration under stop-gradient.

\paragraph{Loss landscapes and basin structure.}
The basin analysis uses linear mode connectivity and permutation alignment
ideas from the neural-network loss-landscape literature
\citep{frankle2020linear,entezari2022role}. Alignment is needed because raw
interpolation between component models can create artificial barriers from
permutation symmetry alone.

\paragraph{Layer-wise objectives.}
Greedy layer-wise pretraining and deep belief networks showed that useful
representations can be learned by protected local objectives
\citep{hinton2006fast}. This paper gives an implicit-EM mechanism for why local
objectives can preserve structure that end-to-end backpropagation does not.

\paragraph{Scope.}
This is not a general theory of deep-network conditioning. The claims are about
EM sites and the corridors through which their responsibility gradients travel.
